\chapter{A Similitude of Numbers}

\lettrine{A}{t} a lively party near the Santa Fe Institute, Murray Gell-Mann’s brainchild, chemists, physicists, and biologists mingled, but where were the mathematicians? 

\begin{figure}[H]
\includegraphics[width=1.0\linewidth]{Illustrations/vign-similitude.png}
\caption{Celiac and Jinyoung discuss turbulence and primes in the kitchen.}
\end{figure}

Naturally, in their dual space of the kitchen. Here Celiac Wu had made her way in search of a quiet corner to retreat from the bustle and bluster from afar. She was contemplating a bit of grape juice when the sink tap suddenly spluttered, projecting a dash of water her way. 

Standing at the sink, looking a little sheepish, was Jinyoung,  visiting the Institute to explore modular forms and the likes. She had turned the tap on too fast whilst rinsing a wine glass, causing the burstful scatter. Celiac laughed softly. "That great mystery of the similitudes\footnote{The principle of similitude, a cornerstone of scaling in the physical sciences, owes its origins to the Latin word \textit{similitudo}, meaning "resemblance." It is the art of drawing inferences about the behavior of systems by identifying dimensionless ratios that remain invariant under changes in scale that reveal relationships that transcend the units of measurement, enabling insights into underlying laws governing systems of different sizes.} and the physics of turbulent flow."

Jinyoung grinned, setting the glass down carefully. "I suppose I should have been more careful. But now that you mention it, I’ve always found turbulence fascinating. It’s chaotic yet structured, much like prime numbers, don’t you think?"

\section*{Turbulence Number Flow}

Celiac leaned against the counter, intrigued. "Exactly. Turbulent flow and laminar flow are governed by dimensionless quantities like the Reynolds number defined as:
\[
\text{Re} = \frac{\rho v L}{\mu},
\]
where \(\rho\) is the fluid density, \(v\) is the velocity of the flow, \(L\) is the characteristic length, and \(\mu\) is the dynamic viscosity. When \(\text{Re}\) is below a critical threshold, the flow remains smooth and laminar. Above it, chaos reigns."

Jinyoung nodded, wiping her hands. 

"So, it’s a balance of forces—viscous forces trying to dampen the flow and inertial forces driving it forward?"

"Precisely," Celiac said, gesturing toward the still-dripping tap. "The beauty of similitude," Celiac continued, "is that it allows us to make sense of how systems scale.  "Similarly, the Froude number describes the influence of gravity on fluid flow, 
\(
\text{Fr} = \frac{v}{\sqrt{gL}},
\)
where \(g\) is the acceleration due to gravity and enables us to compare vastly different sized systems from a tiny laboratory model or a full-scale ship."

Jinyoung nodded thoughtfully. "So, the principle of similitude seeks invariants as dimensionless quantities, that scale naturally across systems. Does it always rely on physical dimensions, or can it apply more abstractly?"

"I could imagine so ... why ... yes, the tap gives us a perfect analogy. When the water flows gently, it’s laminar, like the quiet progression of natural numbers. But turn it up too high, and turbulence emerges and we have eddies, vortices, and chaotic patterns. It’s like prime numbers appearing irregularly among the composites."

Jinyoung’s face lit up. "That’s fascinating. And you know, prime numbers do have their own <<dimensionless>> measure" (bless her she air quoted- it rankled Celiac but not enough)—"the ratio between primes and their logarithmic density:
\[
\pi(x) \sim \frac{x}{\ln(x)},
\]
where \(\pi(x)\) counts the primes less than \(x\). Could we think of this as a Reynolds-like number for the distribution of primes?"

Celiac raised an eyebrow. "That’s an intriguing idea. If turbulence represents chaos but still obeys underlying rules, then the thinning of primes on the number line might also follow some emergent principles. What if there’s a similitude in number theory which is analogous to Reynolds or Froude numbers. One that governs the scaling of primes?"

Jinyoung turned the tap back on, this time letting the water flow gently. "That is rather compelling. In fluid dynamics, similitude lets us compare systems of different scales, just by identifying dimensionless ratios that remain constant. Now thinking now aloud a little pompously, "The thinning of primes as numbers grow, for instance, mirrors the attenuation of physical effects in scaling systems. Could we define a similar principle for numbers? What if we redefine the integer numberline, \(\mathbb{Z}\) as a quotiented space of "dimensionless" 'prime scaled factors'?" For instance:
\begin{itemize}
    \item \(6\) becomes \(\frac{6}{5}\) because \(5\) is the largest prime less than \(6\),
    \item \(9\) becomes \(\frac{9}{7}\) because \(7\) is the largest prime less than \(9\).
\end{itemize}"
"So, by scaling numbers through primes, you’re constructing a 'dimensionless' structure for the number line?"

"Exactly," Celiac replied. "The scaled fractions encode how numbers relate to their prime anchors. This simililtudinal system, \(\mathbb{S}\) highlights the attenuation of primes—how their density decreases as the number line progresses. Celiac grabbed a notebook and began sketching. "Imagine redefining the number line, not as a sequence of integers, but as scaled fractions of primes. Instead of \(\{1, 2, 3, 4, \dots\}\), we construct:
\[
\{1, 2, \frac{3}{2}, \frac{4}{3}, \frac{5}{3}, \frac{6}{5}, \frac{7}{5}, \frac{8}{7}, \frac{9}{7}, \frac{10}{7}, \frac{11}{7}, \frac{12}{11}, \dots\}.
\]

\begin{figure}[H]
\includegraphics[width=1.0\linewidth]{Illustrations/similitude200.png}
\caption{Saw tooth of similitude rationals from n=1 to 200.}
\end{figure}

\begin{itemize}
    \item \textbf{Linear Increase:} Within each interval, \( s(n) \) increases linearly, starting at \( 1 \) for \( n = p_k \) and rising up to \( \frac{p_{k+1} - 1}{p_k} \).
    \item \textbf{Step-Downs:} At each new prime \( p_{k+1} \), \( s(n) \) resets to \( 1 \), resulting in noticeable step-downs in the plot.
    \item \textbf{Prime Density:} The steps are more frequent at smaller \( n \) due to the higher density of primes. As \( n \) increases, primes become sparser, and the intervals widen.
\end{itemize}
\begin{figure}[H]
\includegraphics[width=1.0\linewidth]{Illustrations/similitude10000000.png}
\caption{Similitude rationals across different scales.}
\end{figure}

Here, each integer \(n\) is scaled in the \href{https://colab.research.google.com/drive/1FCwpkh0DpqdNee8hJfEObWrliPObtxUl?usp=sharing}{code} relative to the largest prime less than or equal to it. The ratios could help us study the 'turbulent' distribution of primes—how they cluster and thin out as we move along the number line."

Jinyoung leaned back, a smile spreading across her face. "We could define 'prime density ratios' or even entropy-like measures for composites, using these scaled structures."

The two mathematicians fell silent for a moment. Around them, the party buzzed with idle conversation of the "did you hear what Janet said of Donald" variety. Scientists are only human,

"Maybe," Jinyoung said finally, "the principle of similitude isn’t just a tool for physicists. It’s a means to uncover universal scaling patterns, whether in water or in numbers." Celiac raised her glass in a toast. "Forget our turbulent times, to turbulence and primes."

Jinyoung scribbled one last note before leaving:
\emph{"Similitude isn’t just a principle—it’s an invitation. Whether in nature or numbers, it calls us to seek the universal in the particular, the invariant in the ever-changing."}

% >>>>>> ANNEX C (cut-sheet v2): the similitude power-law fit — block commented out below, pointer left in text <<<<<<
\textit{(The similitude power-law fit --- the full working is set out in Annexe C.)}

% \section*{Postscript}
% 
% The similitude values \( s(n) \) for the given ranges are modeled well by a power law decay of the form:
% \[
% s(n) \sim 1 + a \cdot n^{-b}
% \]
% The fitted parameters \( a \) and \( b \) for each range are as follows:
% \begin{align*}
% \text{Range: } n = 1 \text{ to } 1000 & \quad s(n) \sim 1 + 0.501 \cdot n^{-0.608} \\
% \text{Range: } n = 1001 \text{ to } 10000 & \quad s(n) \sim 1 + 1.828 \cdot n^{-0.860} \\
% \text{Range: } n = 10001 \text{ to } 100000 & \quad s(n) \sim 1 + 2.217 \cdot n^{-0.880} \\
% \text{Range: } n = 100001 \text{ to } 1000000 & \quad s(n) \sim 1 + 2.996 \cdot n^{-0.906}
% \end{align*}
% 
% \begin{itemize}
%     \item exponent \( b \) is converging as \( n \) increases, to around \( b \approx 0.9 \)?
%     \item coefficient \( a \) increases monotonically with \( n \), reflecting the scaling behavior of similitudes across different ranges.
% \end{itemize}
% 
% >>>>>> end of annexed block <<<<<<
\section*{On the Scarcity of Odd Perfect Numbers}

\lettrine{J}{ovannah} leans over the café table, stirring her coffee absentmindedly. Across from her, Celiac flips through a set of scribbled notes on figurative number sequences, barely listening—until Jovannah makes an offhand remark that halts her movements.

\bigskip

\begin{itemize}[leftmargin=1.5cm]

\item[Jovannah:] \textit{(leaning back, setting his spoon down)} "We know the sum of reciprocals of perfect squares converges—\(\sum 1/n^2\) is famously finite. More generally, any figurative number sequence that is quadratic in generator behaves the same way. Even the triangle numbers, dense as they may seem, still lead to a finite sum."

\item[Celiac:] \textit{(raising an eyebrow, setting her notes aside)} "You’re talking about reciprocal density arguments? Sure, perfect numbers are incredibly rare; though, strictly speaking, we still don't know whether an odd perfect number even exists."

\item[Jovannah:] \textit{(grinning slightly)} "And yet, we do know that an \textbf{8-dimensional cuboid number} exists—a structure even more exotic than an odd perfect number. Once we found that 945’s sum of divisors exceeds \(2 \times 945\), all bets were off."

\item[Celiac:] \textit{(scoffing, but intrigued)} "Wait, you’re really going to compare a superabundant number to a perfect one? That’s quite a leap."

\item[Jovannah:] \textit{(tapping the table)} "Not as much as you'd think. The very fact that an integer like 945 defied early assumptions about divisor sums should make us wary of saying \textit{no odd perfect number can exist}. We’ve pushed our searches beyond \(10^{300}\) and found nothing—but absence isn’t disproof. If even one odd perfect number exists, its reciprocal sum, however minuscule, will still contribute to the density conversation."

\item[Celiac:] \textit{(smiling wryly)} "So you’re saying it’s possible that somewhere in the universe of numbers, there's a lone odd perfect number waiting to make a fool out of us?"

\item[Jovannah:] \textit{(chuckling)} "Exactly. We dismissed 945 at first, too and look where that got us."

\end{itemize}

\newpage











